% Generated draft: requires mathematical and source review.
\documentclass[11pt]{article}
\usepackage[margin=25mm]{geometry}
\usepackage{fontspec}
\setmainfont{Latin Modern Roman}
\usepackage{amsmath,amssymb,mathtools,unicode-math}
\setmathfont{Latin Modern Math}
\usepackage{enumitem,needspace,xurl,hyperref}
\setlength{\parindent}{0pt}
\setlength{\parskip}{5pt}
\setlength{\emergencystretch}{3em}
\newcommand{\N}{\mathbb{N}}
\newcommand{\Z}{\mathbb{Z}}
\newcommand{\Q}{\mathbb{Q}}
\newcommand{\R}{\mathbb{R}}
\newcommand{\C}{\mathbb{C}}
\newcommand{\F}{\mathbb{F}}
\newcommand{\A}{\mathbb{A}}
\newcommand{\PP}{\mathbb{P}}
\newcommand{\E}{\mathbb{E}}
\newcommand{\eps}{\varepsilon}
\newcommand{\dd}{\,\mathrm{d}}
\begin{document}
\section*{UnsolvedMath ID: 30001474}
\textbf{Source problem code:} catalogue entry
\subsection{Definitions and mathematical statement}
Use the customary convention that an \emph{open surface} is a connected, noncompact, Hausdorff, second-countable surface without boundary; no orientability assumption is imposed. A smooth structure is understood when discussing a minimal embedding. These conventions are not explicitly specified in the supplied statement and require confirmation against the source context.

An \emph{end} of \(M\) is a Freudenthal end, represented by compatible components outside compact subsets of \(M\). An end neighborhood is a connected neighborhood of that end outside a compact subset. An end is \emph{nonorientable} if every sufficiently small end neighborhood is nonorientable, equivalently if it has no orientable end neighborhood. To make \emph{infinite genus at an end} precise for both orientable and nonorientable surfaces, use the following convention. For a compact connected subsurface \(K\) with \(b(K)\) boundary components, put
\[
\gamma(K)=2-b(K)-\chi(K).
\]
This is twice the usual genus when \(K\) is orientable and is the nonorientable genus when \(K\) is nonorientable. An end has infinite genus if, for every end neighborhood \(U\) and every integer \(n\geq 1\), there is a compact connected subsurface \(K\subset U\) with \(\gamma(K)\geq n\). The supplied record does not explicitly define its end and genus terminology; this convention needs human mathematical review.

A smooth bounded domain \(D\subset\R^3\) is a bounded connected open set with smooth boundary. A \emph{complete proper minimal embedding} \(f\colon M\to D\) is a smooth embedding with vanishing mean curvature whose induced Euclidean metric is complete and for which \(f^{-1}(C)\) is compact in \(M\) for every compact subset \(C\subset D\). Properness is relative to \(D\).

With these conventions, the supplied question asks whether, for every open surface \(M\),
\[
\begin{aligned}
&\exists\,D\subset\R^3\text{ a smooth bounded domain}\;
\exists\,f\colon M\hookrightarrow D
\text{ a complete proper minimal embedding}
\\[-1mm]
&\qquad\Longleftrightarrow\quad
\left\{
\begin{array}{l}
M\text{ has finitely many nonorientable ends},\\
\text{every end of }M\text{ has infinite genus}.
\end{array}
\right.
\end{aligned}
\]
The domain may depend on \(M\), and finitely many nonorientable ends includes zero. This equivalence is posed as a question, not asserted as a theorem.
\subsection{Short English statement}
Does an open surface admit a complete minimal embedding that is proper inside some smooth bounded domain of \(\R^3\) exactly when it has only finitely many nonorientable ends and infinite genus at every end? The surface and end conventions above require source confirmation. The supplied registry labels this problem open; that label has not been independently verified.
\subsection{Sources}
\begin{itemize}
\item Supplied local UnsolvedMath registry record: permanent ID \(30001474\), title \emph{Proper Minimal Embeddings of Open Surfaces}, exact source problem code \texttt{catalogue entry}, in the Oberwolfach Reports Open Problems collection. The proposed year is recorded as 2010; no individual proposer is supplied.
\item Source report: \emph{Progress in Surface Theory} (2010), Oberwolfach Reports, pp.~1225--1318, DOI 10.4171/owr/2010/21. \url{https://doi.org/10.4171/owr/2010/21}. The supplied background describes these reports as documenting mathematical workshops and their open-problem discussions.
\item Supplied literature reference: Meeks--P\'erez--Ros. \url{https://doi.org/10.1215/00127094-2020-0087}. No article title or specific supporting result is supplied.
\item Supplied literature reference: Alarc\'on--L\'opez. \url{https://doi.org/10.1016/j.aim.2012.05.007}. No article title or specific supporting result is supplied.
\item Literature assessment dated exactly \texttt{2026-08-21}: the supplied curation labels the problem open and states that finite-genus cases and proper immersions have substantial results, while the stated classification for proper embeddings of arbitrary open surfaces remains unresolved. This is supplied dated literature triage derived from the OWR curation, not an independently verified assessment; the time-sensitive status and cited claims require independent verification.
\item The registry assigns statement status \texttt{corrected\_verified} and reports that the mathematical content was checked against the official Oberwolfach report, with notation, context, and question form normalized without changing the source-level claim. This is supplied verification metadata, not a verification performed for this draft.
\end{itemize}
\end{document}
